\chapter{Literature Review}

In this chapter, we review the relevant literature on multi-currency 
portfolio optimization. We discuss the contributions and 
methodologies in the field, highlighting how our work builds 
upon and extends existing research. 

The core of our analysis builds on two recent studies that 
apply regularization techniques to international asset allocation. 
Burkhardt and Ulrych (2023) introduce a joint optimization approach 
within a mean-variance framework. By using L1 and L2 regularization, they show that jointly optimizing assets and 
currencies leads to better out-of-sample performance than the standard 
industry practice of using a separate currency overlay. 

Extending this framework, Lucescu and Ulrych (2026) consider an 
Expected Shortfall optimization framework to better manage downside 
tail risk. To fix the numerical instability common in Expected Shortfall minimization, 
they apply similar regularization methods. Their findings reveal a 
dimensional trade-off. Joint optimization works best in small asset universes, 
but the separate currency overlay approach becomes more effective for larger 
universes due to its built-in dimensionality reduction. 

Reproducing and comparing these two baseline frameworks serves as the starting point 
for the extensions proposed in this thesis.
